\documentclass[10pt,a4paper]{amsart}
\usepackage[margin=19mm]{geometry}
\usepackage{amsmath,amssymb,mathtools}
\usepackage[scr=rsfs,cal=euler]{mathalfa}
\usepackage{xfrac}
\usepackage[unicode,hidelinks,bookmarks=false]{hyperref}
\newcommand{\ie}{i.e.\ }
\newcommand{\cf}{cf.\ }
\newcommand{\resp}{respectively\ }
\newcommand{\Subsection}{\subsection}
\providecommand{\nobreakdash}{\nobreak\hbox{-}}
\setlength{\emergencystretch}{2em}
\multlinegap0pt
\usepackage{bm}
\usepackage{bbm}
\usepackage{dsfont}
\usepackage{xspace}
\usepackage{cancel}
\usepackage{mathtools}
\usepackage{amsthm}
\makeatletter
\@ifundefined{lemma}{\newtheorem{lemma}{Lemma}}{}
\makeatother
\makeatletter
\def\E{\mathbb E}
\def\R{\mathbb R}
\def\d{\delta}
\renewcommand{\ge}{\geqslant}
\renewcommand{\le}{\leqslant}
\expandafter\ifx\csname proof\endcsname\relax
\newenvironment{proof}
   %   {\medskip\noindent{\bf Proof:}\hspace{1mm}}
\fi
\expandafter\ifx\csname proofof\endcsname\relax
\newenvironment{proofof}[1]
%       {\medskip\noindent{\bf Proof of #1:}\hspace{1mm}}
%       {\hfill$\Box$\medskip}
% \input{epsf}
\fi
\makeatother
\title[A fractal-like configuration of point-line pairs for the min]{A fractal-like configuration of point-line pairs for the minimal distance problem}
\author{Alexander Logunov, Dmitrii Zakharov}
\date{Source: arXiv:2511.10509v1 (CC BY 4.0)}
\begin{document}
\maketitle

\noindent\emph{Source and licence.} A fractal-like configuration of point-line pairs for the minimal distance problem. Version arXiv:2511.10509v1 (CC BY 4.0), distributed under the Creative Commons Attribution 4.0 International licence, as recorded in the arXiv licence metadata for this exact version and shown on the abstract page https://arxiv.org/abs/2511.10509v1. Full rights notice: licenses/arxiv-2511.10509-CC-BY-4.0.txt.\par\smallskip

\noindent\textit{Editorial scope.} Counted unit: the Lemma whose source label is \texttt{lem1} in the article (source0.tex lines 174--176, source label \texttt{lem1}), delivered together with the article's own complete original proof of it (source0.tex lines 196--232, one \texttt{proof} environment of 37 lines). No mathematics is written, repaired or re-derived: statement and proof are the article's own text line for line. The proof is detached from the statement in the source: the article states the result at lines 174--176 and prints its proof at lines 196--232, and the proof environment's own header names the statement, so the association is the article's own. Required context retained uncounted: none. Every definition, display and label of those lines is retained unchanged. Omitted from this package and not redistributed: 1--173, 177--195, 233--302. Nothing inside a retained excerpt is omitted or altered: every retained body is delivered exactly as the article's lines are written, so no omission receipt of any kind accompanies this package. Citations: the retained excerpts carry no citation command at all, so no bibliography entry of the article is transcribed and no reference is redistributed. Reference adaptation: none was needed and none was made; every reference inside the delivered unit resolves inside it, so no unresolved reference or citation is produced. Printed slips kept literal: no printed word, symbol, hypothesis or number of the counted statement or of its proof was altered. Layout and class: the article's own class is \texttt{amsart} with packages including tikz, amsfonts, amsmath, latexsym, color, epsfig, hyperref, enumitem, amssymb, bbm, scalerel, stackengine, esint; the wrapper is the lane's pinned \texttt{amsart} editorial wrapper at 10pt on A4, which restates the theorem-like environments the retained text needs and adds the article's own macro definitions that the retained text needs (\texttt{\textbackslash E}, \texttt{\textbackslash R}, \texttt{\textbackslash d}, \texttt{\textbackslash ge}, \texttt{\textbackslash le}), with extra wrapper packages (bm, bbm, dsfont, xspace, cancel, mathtools). The delivered numbering is the unit's own. Assets: no retained span contains a figure environment, a graphics-inclusion command, a PDF-page inclusion, a table or a TikZ picture; no image file is copied into this package, the package declares no asset dependency and no image file is redistributed with it. Licence: version arXiv:2511.10509v1 of this article is distributed under the Creative Commons Attribution 4.0 International licence, as recorded in the arXiv licence metadata for this exact version and shown on the abstract page https://arxiv.org/abs/2511.10509v1. A full rights notice is delivered at licenses/arxiv-2511.10509-CC-BY-4.0.txt. Characters: every retained excerpt is pure ASCII, so the delivered engine, XeLaTeX with its default Latin Modern text font, reports no missing character.
\begin{lemma}\label{lem1}
    There exists a constant $c_1>0$ such that $n(\d) \ge c_1 \d^{-1} \log 1/\d$ for all $\d\in (0,1)$.
\end{lemma}

\begin{proof}[Proof of Lemma \ref{lem1}]
    If is enough to prove the lemma for all sufficiently small $\d \le\d_0$ (for $\d>\d_0$ we can just choose $c_1$ small enough). For $p\in \R^2$ and $\theta \in [-1, 1]$ let $T(p, \theta) = p + (1, \theta)\R + \{0\} \times [-\d, \d]$, i.e. $T(p, \theta)$ is a strip of vertical width $2\d$ through $p$ at slope $\theta$.
    
    Let $N = [0.1 \d^{-1} (\log 1/\d)]$. 
    Let $p_1, \ldots, p_N \in [0,1]^2$ be a uniformly random and mutually independent collection of points and let $P = \{p_1, \ldots, p_N\}$. 
    For each point $p_i$, it suffices to show that the probability that there exists an angle $\theta\in [-1,1]$ such that $T(p_i, \theta) \cap P = \{p_i\}$ is at least $1/2$ (for $\d\le \d_0$). Indeed, by linearity of expectation this would give a collection of $N/2$ points among $p_1, \ldots, p_N$ so that for every $p_i$ in the collection there is an angle $\theta$ such that $T(p_i, \theta)\cap P =\{p_i\}$. In particular, this gives $n(\d) \ge N/2 \ge c \d^{-1}\log 1/\d$.
    
    Fix an index $i$. 
    Let $Q_i = p_i + [- \frac{1}{4\sqrt{N}}, \frac{1}{4\sqrt{N}}]^2$. 
    Note that for $i'\neq i$ we have $p_{i'}\in Q_i$ with probability at most $|Q_i|\le \frac{1}{4N}$. So by the union bound we have
    \begin{equation*}\label{eq:Q}
    \Pr[Q_i \cap P \neq \{p_i\}] \le (N-1) |Q_i| \le 1/4.    
    \end{equation*}
    So with probability at least $3/4$ there are no points inside $Q_i$ besides $p_i$. So it suffices to show that with probability at least $3/4$ there exists a slope $\theta$ such that $(T(p_i, \theta) \setminus Q_i)\cap P = \emptyset$. 

    Let $P' = P\setminus\{p_i\}$.
    Denote $M = [\frac{1}{8\d\sqrt{N}}]$ and let $\theta_j = 8j \d \sqrt{N}$, for $j=1, \ldots, M$. Note that for any $j\neq j' \in \{1, \ldots, M\}$ we have $T(p_i, \theta_j) \cap T(p_i, \theta_{j'})\subset  Q_i$. We now estimate the probability that $P' \cap (T(p_i, \theta_j) \setminus Q_i) = \emptyset$ for some $j \in \{1, \ldots, M\}$. 
    For a region $A \subset \R^2$ the probability that $A\cap P' = \emptyset$ is equal to $(1 - |A \cap [0,1]^2|)^{N-1}$. By mutual independence of the collection $\{p_1 \ldots, p_N\}$, this also holds if $A$ is a random set depending only on $p_i$. Since $T(p_i,\theta)$ has vertical width $2\d$, we have $|T(p_i,\theta) \cap [0,1]^2| \le 2\d$. So we can choose pairwise disjoint sets $A_j\subset [0,1]^2$ of area $2\d$ so that $A_j$ contains $T(p_i, \theta_j) \cap [0,1]^2 \setminus Q_i$.
    
    
    %We can take $M \ge \frac{1}{10\d \sqrt{N}}$. 
    
    Let $Z_j = 1_{A_j \cap P' = \emptyset}$ be the indicator of the event that $A_j \cap P' = \emptyset$. Then we have 
    \[
    \E Z_j = (1-2\d)^{N-1} \ge \exp(- 2\d (N-1)) \ge \exp(-2\d \times 0.1\d^{-1} \log (1/\d)) \ge \d^{0.2}.
    \]
    Denote $q = (1-2\d)^{N-1}$. So the sum $Z= \sum_{j=1}^M Z_j$ has expectation $\E Z =M q$. Since $M = [\frac{1}{8\d\sqrt{N}}] \sim (\d \log(1/\d))^{-1/2} \gg \d^{-0.2}$, we have $Mq > 16$ for $\d \le \d_0$. We note that the variables $Z_j$ are negatively correlated: for $j\neq j'$ the product $Z_{j} Z_{j'}$ is the indicator of the event that $P' \cap (A_{j}\cup A_{j'}) = \emptyset$ and so $\E Z_j Z_{j'} = (1-4\d)^{N-1} \le q^2 = (\E Z_j) (\E Z_{j'})$. For $j=j'$ we have $\E Z_jZ_{j'} = \E Z_j = q$. 
    We use this to estimate the variance of $Z$:
    \begin{align*}
    \operatorname{Var}Z = \E Z^2 - (\E Z)^2= \sum_{j, j'=1}^M \left(\E Z_j Z_{j'}- (\E Z_j)(\E Z_{j'}) \right)\le  M q.
    \end{align*}
    By Chebyshev's inequality,
    \[
    \Pr[ |Z -\E Z| \ge \E Z/2] \le 4\frac{\operatorname{Var}Z}{(\E Z)^2} \le \frac{4}{M q} \le 1/4.
    \]
    We conclude that for every $i$, the probability that there exists $\theta$ with $T(p_i, \theta)\cap P=\{p_i\}$ is at least $1/2$ which in turn implies that $n(\d) \ge N/2 \ge c\d^{-1}\log 1/\d$, as desired.
\end{proof}

\end{document}
